\DocumentMetadata{
  pdfversion=2.0,
  pdfstandard=ua-2,
  lang=en-US,
  tagging=on,
  tagging-setup={math/setup=mathml-SE,tabsorder=structure}
}
\documentclass[11pt,letterpaper]{article}
\usepackage[margin=0.68in,includefoot]{geometry}
\usepackage{newcomputermodern}
\usepackage{amsmath,amscd,mathtools}
\usepackage{tikz,adjustbox}
\providecommand{\square}{\mdlgwhtsquare}
\usepackage[hidelinks]{hyperref}
\hypersetup{
  pdftitle={Numerical Analysis PhD Exam, May 2017, Part 2},
  pdfauthor={University of Florida Department of Mathematics},
  pdfsubject={Graduate mathematics examination},
  pdfdisplaydoctitle=true,
  bookmarksopen=true
}
\setcounter{secnumdepth}{0}
\setlength{\parindent}{0pt}
\setlength{\parskip}{0.28em}
\setlength{\emergencystretch}{3em}
\setlength{\leftmargini}{2.15em}
\setlength{\leftmarginii}{2.65em}
\setlength{\leftmarginiii}{3em}
\pagestyle{empty}
\raggedbottom
\allowdisplaybreaks
\NewTaggingSocketPlug{block/list/label}{examproblem}{%
  \tagstructbegin{tag=Lbl}\tagstructend%
  \tagstructbegin{tag=\LBody}%
  \tagstructbegin{tag=\ProblemHeadingTag,title-o={\ProblemHeadingText}}%
  \tagmcbegin{tag=\ProblemHeadingTag,actualtext={\ProblemHeadingText}}%
  #2%
  \tagmcend%
  \tagstructend%
}
\newenvironment{examproblems}[2]{%
  \def\ProblemHeadingTag{#2}%
  \AssignTaggingSocketPlug{block/list/label}{examproblem}%
  \begin{enumerate}%
  \setcounter{enumi}{#1}%
  \renewcommand{\labelenumi}{\textbf{\theenumi.}}%
  \setlength{\topsep}{0.35em}%
  \setlength{\partopsep}{0pt}%
  \setlength{\itemsep}{0.48em}%
  \setlength{\parsep}{0.18em}%
}{\end{enumerate}}
\newenvironment{examsubparts}{%
  \AssignTaggingSocketPlug{block/list/label}{default}%
  \begin{enumerate}%
  \setlength{\topsep}{0.18em}%
  \setlength{\partopsep}{0pt}%
  \setlength{\itemsep}{0.16em}%
  \setlength{\parsep}{0.1em}%
}{\end{enumerate}}
\newenvironment{examinstructions}{%
  \par\begingroup\itshape
}{%
  \par\endgroup\vspace{0.18em}
}
\makeatletter
\renewcommand\subsection{\@startsection{subsection}{2}{0pt}%
  {-1.25ex plus -.25ex minus -.1ex}%
  {0.55ex plus .1ex}%
  {\normalfont\large\bfseries}}
\renewcommand{\@maketitle}{%
  \newpage\null\vskip -1em%
  {\footnotesize\bfseries
    \href{https://www.ufl.edu/}{University of Florida}, \href{https://math.ufl.edu/}{Department of Mathematics}\par}%
  \vskip -0.5em%
  \tagstructbegin{tag=H1,title={Numerical Analysis PhD Exam, May 2017, Part 2}}%
  \tagpdfparaOff%
  \tagmcbegin{tag=H1}%
    {\LARGE\bfseries\href{https://gma.math.ufl.edu/exams/numerical-analysis-phd/}{Numerical Analysis PhD Exam}, May 2017, Part 2\par}%
  \tagmcend%
  \tagpdfparaOn%
  \tagstructend%
  \vskip 0.25em%
  \tagmcbegin{artifact}%
    \hrule height 1pt%
  \tagmcend%
  \vskip 1em%
}
\makeatother
\title{Numerical Analysis PhD Exam, May 2017, Part 2}
\author{}
\date{}
\begin{document}
\maketitle
\thispagestyle{empty}
\begin{examinstructions}
Do all five (5) problems.
\end{examinstructions}
\begin{examproblems}{0}{H2}
\def\ProblemHeadingText{Problem 1.}
\item
With
\[
\phi(w,t)=a f(w+ch,t+bh),
\]
 find the values of the parameters \(a,b,c\) so that the resulting one-step method
\[
w_0=\alpha
\]

\[
w_{i+1}=w_i+h\phi(w_i,t_i)
\]
 has local truncation error \(O(h^2)\).
\def\ProblemHeadingText{Problem 2.}
\item
Hint: the error term in the one unit Trapezoid rule is \(-\frac{h^3}{12}f''(\eta)\) and in the one unit Simpson’s rule is \(-\frac{h^5}{90}f^{(4)}(\eta)\).

\par
Let~\(S\) be the cubic spline given by
\[
S(x)=(x+1)^3 \text{ for } x\in[-1,0]
\]

\[
S(x)=(1-x)^3 \text{ for } x\in[0,1].
\]
\begin{examsubparts}
\item[\textnormal{(a)}]
Estimate the error of the composite trapezoidal rule applied to \(\int_{-1}^1 S(x)\,dx\), when \([-1,1]\) is divided into~\(n\) subintervals of equal length \(h=2/n\) and~\(n\) is even (and so 0 is a node).
\item[\textnormal{(b)}]
Estimate the error of the composite Simpson’s rule applied to \(\int_{-1}^1 S(x)\,dx\), when \([-1,1]\) is divided into~\(n\) subintervals of equal length \(h=2/n\) and~\(n\) is divisible by 4 (and so 0 is a node).
\end{examsubparts}
\def\ProblemHeadingText{Problem 3.}
\item
This problem has two parts.
\begin{examsubparts}
\item[\textnormal{(a)}]
If \(f\in C^1[a,b]\) and \(a\leq x_0<\cdots<x_n\leq b\) and \(H,G\) are degree at most \(2n+1\) polynomials with \(G(x_i)=H(x_i)=f(x_i)\) and \(G'(x_i)=H'(x_i)=f'(x_i)\) for all~\(i\), then \(G=H\).
\item[\textnormal{(b)}]
If \(\varphi_0,\varphi_1,\ldots,\varphi_n\) are polynomials with \(\varphi_i\) of degree~\(i\), then the set of \(\varphi_i\) is linearly independent.
\end{examsubparts}
\def\ProblemHeadingText{Problem 4.}
\item
Consider the inner product on \(C[0,\infty)\) given by
\[
\langle f,g\rangle=\int_0^\infty f(x)g(x)e^{-x}\,dx.
\]
\begin{examsubparts}
\item[\textnormal{(a)}]
Starting with the basis \(\{1,t,t^2\}\) for \(\mathcal P_2[0,\infty)\), find three orthonormal polynomials \(\phi_0,\phi_1,\phi_2\) on \([0,\infty)\) with respect to the inner product and the degree of \(\phi_n\) is equal to~\(n\). Hint: \(\int_0^\infty t^m e^{-t}\,dt=m!\).
\item[\textnormal{(b)}]
Find the equations satisfied by the values of \(w_1,w_2,t_1\) and \(t_2\) which yield the weighted Gaussian Quadrature formula
\[
\int_0^\infty f(t)e^{-t}\,dt=w_1f(t_1)+w_2f(t_2)
\]
 with degree of precision 3.
\end{examsubparts}
\def\ProblemHeadingText{Problem 5.}
\item
This problem has two parts.
\begin{examsubparts}
\item[\textnormal{(a)}]
Assume \(g\in C^2[a,b]\) with \(g([a,b])\subset[a,b]\) and fixed point \(p\in(a,b)\). Assume that \(g'(p)=0\). Show that for any \(x\in[a,b]\) with \(x\ne p\),
\[
\frac{|g(x)-p|}{|x-p|^2}\leq M,
\]
 where \(M=\max\{|g''(z)|:z\in[a,b]\}/2\).
\item[\textnormal{(b)}]
Let \(g(x)=x-\tan(x)\). Find a fixed point~\(p\) of~\(g\) with \(g'(p)=0\) and give an explicit \([a,b]\) where you prove that for all \(x\in[a,b]\), \(g^n(x)\to p\).
\end{examsubparts}
\end{examproblems}
\end{document}
